\section{Rangkuman}
\label{sec:rangkuman-bab3}

Pada bab ini telah dibahas:

\begin{itemize}
    \item Bahasa pemrograman dibedakan menurut tingkat abstraksinya:
          bahasa mesin, bahasa assembly, dan bahasa tingkat tinggi.
          Bahasa assembly berada tepat di atas bahasa mesin dan
          bersifat \eng{one-to-one} terhadapnya.

    \item Bahasa assembly tetap dipakai ketika dibutuhkan kendali
          langsung atas perangkat keras, ukuran program yang kecil,
          serta pemahaman cara kerja prosesor pada level instruksi.

    \item Bahasa assembly 8086 mengenal bilangan biner, oktal,
          desimal, dan heksadesimal. Setiap basis mempunyai aturan
          digit dan bobot posisi yang berbeda.

    \item Konversi dapat dilakukan dengan menjumlahkan bobot posisi,
          membagi berulang dengan basis tujuan, atau mengelompokkan
          bit menjadi tiga bit untuk oktal dan empat bit untuk
          heksadesimal.

    \item Assembler menerima literal desimal tanpa akhiran atau
          dengan akhiran \code{D}, biner dengan akhiran \code{B}, dan
          heksadesimal dengan akhiran \code{H}. Heksadesimal yang
          diawali huruf A--F harus diberi nol di depan.

    \item Bilangan tidak bertanda memakai seluruh bit sebagai
          besaran. Bilangan bertanda memakai bit paling kiri sebagai
          bit tanda dan sisanya mengikuti aturan
          \eng{two's complement}.

    \item Pola bit yang sama dapat ditafsirkan berbeda. \hex{F3}
          berarti 243 sebagai bilangan tidak bertanda dan $-13$
          sebagai bilangan bertanda 8 bit.

    \item Karakter disimpan sebagai kode ASCII satu byte. \code{'A'}
          disimpan sebagai \hex{41}, sehingga data karakter dan data
          bilangan dapat diperiksa dengan cara yang sama di memori.
\end{itemize}

\begin{checklistbox}
\vspace{1.1em}
Tandai setiap butir yang sudah Anda kuasai.

\begin{itemize}
    \item[\checkbox] Saya dapat menjelaskan tiga tingkat bahasa
          pemrograman dan kedudukan assembly.
    \item[\checkbox] Saya dapat mengonversi bilangan biner, oktal,
          desimal, dan heksadesimal.
    \item[\checkbox] Saya dapat menuliskan literal bilangan yang
          diterima assembler, termasuk nol di depan heksadesimal.
    \item[\checkbox] Saya dapat membedakan bilangan bertanda dan
          tidak bertanda.
    \item[\checkbox] Saya dapat membentuk \eng{two's complement} dan
          menafsirkan pola bit seperti \hex{F3}.
    \item[\checkbox] Saya dapat menjelaskan hubungan kode ASCII dengan
          isi memori.
    \item[\checkbox] Saya dapat membaca hasil \file{d ds:0} pada DEBUG
          dan menjelaskan setiap byte.
\end{itemize}

\vspace{0.4em}
Bila masih ada butir yang belum bertanda, ulangi bagian yang
bersesuaian sebelum melanjutkan ke bab berikutnya.
\end{checklistbox}

\begin{infobox}
\vspace{1.1em}
\textbf{Rujukan bab.} Materi pada bab ini diringkas dari
\cite{gunawan2019}, Bagian 1, hlm.~10--11 dan Bagian 14.1,
hlm.~127--128, serta \cite{sto2001}, BAB I, hlm.~6--8, BAB V,
hlm.~23--24, dan BAB XV, hlm.~78.
\end{infobox}
